\documentclass[10pt,letterpaper]{article}

%
% Packages
%
\usepackage[margin=1in]{geometry}
\usepackage{xcolor}
\usepackage[colorlinks=true,linkcolor=black,citecolor=blue,urlcolor=blue]{hyperref}
\usepackage{enumitem}
\usepackage{titlesec}
\usepackage{etoolbox}

%
% Bibliography (reuses the website's shared .bib files).
%
\usepackage[backend=biber,style=numeric,sorting=none,giveninits=true,
            maxbibnames=99,mincrossrefs=99,minxrefs=99,
            doi=false,isbn=false,url=false]{biblatex}
\addbibresource{../assets/pubs/pub.bib}
\addbibresource{../assets/pubs/conf.bib}
\addbibresource{extra.bib}

% Strip co-first-author asterisks from names so biber parses them cleanly.
\DeclareSourcemap{
  \maps[datatype=bibtex]{
    \map{\step[fieldsource=author, match=\regexp{\*}, replace={}]}
  }
}

% Bold the author's own name ("Insu Yun") in the reference list.
\def\myfamilyname{Yun}
\def\mygivenname{Insu}
\DeclareNameFormat{author}{%
  \ifboolexpr{ test {\ifdefstring{\namepartfamily}{\myfamilyname}}
           and test {\ifdefstring{\namepartgiven}{\mygivenname}} }
    {\mkbibbold{\usebibmacro{name:family-given}
        {\namepartfamily}{\namepartgiveni}{\namepartprefix}{\namepartsuffix}}}
    {\usebibmacro{name:family-given}
        {\namepartfamily}{\namepartgiveni}{\namepartprefix}{\namepartsuffix}}%
  \usebibmacro{name:andothers}%
}
\renewcommand*{\bibfont}{\footnotesize}
\setlength{\bibitemsep}{2pt}

%
% Compact spacing to help fit two pages
%
\setlist{noitemsep,topsep=2pt,leftmargin=1.5em}
\titlespacing*{\section}{0pt}{1.2ex}{0.6ex}
\titlespacing*{\subsection}{0pt}{0.8ex}{0.4ex}
\titleformat{\section}{\normalfont\large\bfseries}{\thesection}{0.6em}{}
\titleformat{\subsection}{\normalfont\bfseries}{\thesubsection}{0.6em}{}
\setlength{\parskip}{0pt}
\setlength{\parindent}{1.2em}

%
% \PP{Lead-in} -- a bold run-in paragraph header followed by a period.
%
\newcommand{\PP}[1]{\par\smallskip\noindent\textbf{#1.}\ }

%
% Custom header: title, then name + email on one line, then affiliation.
%
\newcommand{\makeheader}{%
  \begin{center}
    {\LARGE\bfseries Research Overview}\\[2.5ex]
    {\large Insu Yun}\quad\textbar\quad\href{mailto:insuyun@kaist.ac.kr}{insuyun@kaist.ac.kr}\\[0.2ex]
    Associate Professor at KAIST School of Electrical Engineering
  \end{center}
  \vspace{1ex}
}

\begin{document}
\makeheader

% NOTE: Target length is 2 pages. If it overflows, tighten geometry margins,
% trim the prose below, or reduce \parskip.

\section{Overview}
My research seeks to make software security \emph{scientific and scalable},
replacing the manual, intuition-driven process of finding and fixing
vulnerabilities with systematic, automated, and ultimately autonomous methods.
My background as a competitive hacker informs this agenda: I design and build
systems that discover and remediate vulnerabilities in real-world software, and I
draw research questions from the practical difficulty of doing so. This statement
outlines three directions: (1)~automating the discovery of software
vulnerabilities, (2)~advancing automation toward autonomy through AI, and
(3)~applying rigorous security analysis to deployed systems of societal
importance, with implications that extend into public policy.

\section{Automating the Discovery of Software Vulnerabilities}
\PP{Motivation} Modern software has grown too large and changes too rapidly for
manual auditing to keep pace. A central question of my work is how to enable
automated systems to search for vulnerabilities with the depth and adaptivity of
an expert analyst, while retaining the scale that only automation provides.

\PP{QSYM} My most influential contribution in this area is QSYM~\cite{yun:qsym}, a
concolic execution engine designed to enable hybrid fuzzing. Fuzzing is efficient
but shallow, whereas symbolic execution is precise but prohibitively expensive;
QSYM reconciles this tension through a fast, tightly integrated symbolic executor
that allows fuzzers to reach program states previously considered unreachable at
scale. The work helped establish hybrid fuzzing as a practical, mainstream
technique. It received a Distinguished Paper Award at USENIX Security 2018 and
was widely adopted as a building block within the community, and the broader line
of work was subsequently recognized with the Frontiers of Science Award at the
inaugural International Congress of Basic Science (ICBS). QSYM exemplifies the
kind of research I pursue: ideas that are conceptually clean yet directly
useful~in~practice.

\PP{Pushing Automation into Hard Targets} Building on this foundation, I have
developed a broad program spanning diverse domains. I have extended automated
testing to targets long regarded as too complex or stateful to fuzz effectively,
including frameworks for systematically evaluating secure memory
allocators~\cite{yun:hardsheap}, WebAssembly runtimes~\cite{park:rgfuzz}, and
real-time operating-system kernels~\cite{lee:rtcon}. A recurring difficulty is
analyzing software whose source code is unavailable; my group therefore develops
techniques to analyze and query decompiled and binary-only code in order to find
vulnerabilities in commercial off-the-shelf software~\cite{han:queryx}, and has
identified design-level weaknesses in widely deployed systems software, ranging
from cross-cache attacks on the Linux kernel~\cite{kim:crossx} to structural
flaws in the Windows file-system security model~\cite{kim:jenga}. I have further
applied these methods to critical yet under-examined infrastructure: a sustained
line of work analyzes cellular baseband software---the protocol stacks present in
every mobile device---by comparing implementations against their
specifications~\cite{kim:basespec,kim:basecomp}, testing them dynamically over the
air~\cite{park:doltest,hoang:llfuzz,park:otabase}, and emulating their internal
protocol states~\cite{jeong:firmstate}. This breadth also reaches beyond
conventional software: I have synthesized attacks against the automated market
makers that move real value in decentralized finance~\cite{han:cpmm}.

\PP{Broader impact} These contributions extend beyond publications. My engagement
with automated security is long-standing and formative: as a Ph.D. student, I
competed as a finalist in DARPA's Cyber Grand Challenge (CGC), the first competition in
which fully autonomous systems discovered and patched software vulnerabilities
without human intervention. That experience directly motivated the development of
QSYM and, in time, my participation in AIxCC, which I discuss later. Beyond such
competitions, the methods my group develops have repeatedly uncovered and
responsibly disclosed serious vulnerabilities in widely used commercial products,
demonstrating that they operate on real, deployed systems rather than on
benchmarks alone. This coupling of offensive practice with rigorous research is
central to how I work: it keeps my research anchored in consequential, real-world
problems.

\section{Toward Autonomous Cyber Defense with AI}
\PP{Motivation} Existing automation still depends on human involvement: people
decide what to test, interpret the results, and produce the fixes. The next step
is autonomy---systems capable of independently discovering, diagnosing, and
repairing vulnerabilities. Recent progress in large language models (LLMs) makes
this goal newly attainable and, I expect, will substantially change how software
is secured.

\PP{AIxCC} DARPA's AI Cyber Challenge (AIxCC) posed this problem at scale:
whether an autonomous system could secure large, unfamiliar codebases by
discovering, proving, and repairing vulnerabilities without human intervention.
As part of Team Atlanta, I helped develop ATLANTIS, an autonomous
cyber-reasoning system for the competition~\cite{kim:atlantis}. ATLANTIS
integrates my group's expertise in program analysis and fuzzing with the
reasoning capabilities of LLMs into an end-to-end pipeline that discovers
vulnerabilities, identifies their root causes, and generates and validates
patches; its automated repair component, PatchIsland, orchestrates multiple LLM
agents to remediate vulnerabilities at scale and without human
intervention~\cite{kim:patchisland}. Classical program analysis and fuzzing
provide the rigor that grounds the system, while LLMs contribute breadth and
flexible reasoning. ATLANTIS placed first in AIxCC, demonstrating that
autonomous, AI-driven security is practically achievable.

\PP{Broader impact} I believe this work has helped shape the emerging area of
AI-driven vulnerability discovery, and my group continues to advance it. One
direction is \emph{agentic fuzzing}, in which agents---seeded by historical
vulnerabilities---reason about a defect's root cause, hypothesize where comparable
defects may occur elsewhere in a codebase, and validate each hypothesis by
generating and executing proof-of-concept code~\cite{park:afuzz}. More broadly, I
regard the integration of classical program analysis with learned reasoning,
toward systems that autonomously discover and repair vulnerabilities, as one of
the central problems in security for the coming decade.

\section{Strengthening the Security of South Korea}
\PP{Why this matters} Much security research is conducted on abstract,
well-defined problems, yet the failures that cause real harm are often bound up
with how a given society designs and regulates its systems. Problems rooted in a
specific country or context are frequently overlooked by the broader research
community, even though addressing them can prevent genuine harm and reshape
practice in ways that abstract problems rarely do. I therefore consider it an
important responsibility of security researchers to engage with these concrete,
high-impact problems, and my work on South Korea's security ecosystem illustrates
this view.

\PP{Korea Security Applications} In South Korea, users have long been required to
install dedicated security programs---collectively, Korea Security Applications
(KSAs)---to access financial services, software widely assumed, by mandate and
convention rather than evidence, to improve their safety. Because this uniquely
Korean arrangement was largely taken for granted, its actual security
implications had never been rigorously examined. In our study, \emph{Too Much of a
Good Thing: (In-)Security of Mandatory Security Software for Financial Services in
South Korea}~\cite{yun:ksa}, we analyzed this ecosystem systematically and showed
empirically that these KSAs can in fact \emph{undermine} the security they are
intended to provide, enlarging the attack surface and introducing new risks. The result overturns a long-standing national
assumption and illustrates the value of applying the methods of vulnerability
research not to an individual program but to a deployed policy affecting an entire
country.

\PP{Broader impact} This work has had influence well beyond the research
community. Most notably, the President of South Korea has since directed that
these mandatory security applications be phased out---a concrete national policy
outcome that our research helped bring about, and a demonstration that rigorous
security research can effect tangible change. More broadly, I have contributed to
reforming how vulnerabilities are handled in Korea: I helped revise the national bug-bounty program operated by the Korea
Internet \& Security Agency (KISA), replacing a policy of indefinite
non-disclosure with one that discloses vulnerabilities 120 days after a patch
becomes available, and I helped introduce Coordinated Vulnerability Disclosure
(CVD) and Vulnerability Disclosure Policy (VDP) frameworks domestically, bringing
national practice into alignment with international norms. On the basis of this
work, I was appointed to the security special committee under the national AI
strategy committee, where I contribute to national security policy. This
trajectory reflects what I believe security research can ultimately achieve:
addressing the concrete problems a society faces and producing evidence
sufficient to change not only systems but the decisions made about them.

\printbibliography[title={References}]

\end{document}
